\section{Introduction}

Retrieval-Augmented Generation systems answer a query by retrieving
documents from a corpus and conditioning a language model's response on
their contents. This architecture is now standard in production
assistants, but it inherits a quiet assumption: that every document in the
corpus is trustworthy. When that assumption fails, the result is
\textit{knowledge poisoning}, the deliberate insertion of a document
designed to manipulate downstream generation. This is not a hypothetical
concern; indirect prompt injection through externally sourced content has
been demonstrated against production LLM-integrated applications
\cite{yi2023benchmarking}, and dedicated RAG poisoning attacks that require
injecting only a single malicious passage per targeted query have since
been shown to be practical at scale \cite{zhang2025practical}.

A Poison Detection System (PDS) sits between the retriever and the
generator and screens each document before it is trusted. The individual
signals available for this task, keyword and instruction-pattern
matching, embedding-based anomaly detection, and language-model fluency
scoring, are each described in prior work. Our purpose is not to introduce
a fourth signal, but to ask: when these signals are combined, which ones
actually matter, why does a naive combination underperform, and where does
the resulting system still fail, including when it is tested on documents
unlike those it was built on.

\subsection{Contributions}
\begin{enumerate}
  \item A working three-detector PDS pipeline (pattern, embedding,
    perplexity) evaluated on a corpus of 55 real open-source README
    paragraphs and 60 constructed poisoned documents across four attack
    categories, plus 10 adversarial and 12 cross-domain documents.
  \item A diagnosed and corrected fusion failure: equal-weight averaging
    suppresses the most informative detector purely because of a
    raw-score-range mismatch, fixed by calibration and independently
    confirmed by a learned fusion model (F1 0.79 $\rightarrow$ 0.97 and
    0.98).
  \item A staged cascade reaching the same F1 as calibrated fusion with
    70.5\% less detector latency.
  \item An adversarial evaluation that names a specific weakness
    (boilerplate anchoring), a documented calibration-fragility analysis,
    and two honest negative results: segment-level scoring does not fix
    anchoring, and the embedding detector produced false positives on
    every benign document from a different domain.
\end{enumerate}
